\subsection{Penjumlahan Matriks}

Saat menjumlahkan matriks, semua matriks dalam penjumlahan harus berukuran sama.
Sebagai contoh,
\begin{equation*}
\leftB
\begin{array}{rr}
1 & 2 \\
3 & 4 \\
5 & 2
\end{array}
\rightB 
\end{equation*}
dan
\begin{equation*}
\leftB
\begin{array}{rrr}
-1 & 4 & 8\\
2 & 8 & 5
\end{array}
\rightB 
\end{equation*}
tidak dapat dijumlahkan karena yang satu berukuran $3 \times 2$, sedangkan yang lain berukuran $2 \times 3$.

Namun, penjumlahan
\begin{equation*}
\leftB
\begin{array}{rrr}
4 & 6 & 3\\
5 & 0 & 4\\
11 & -2 & 3
\end{array}
\rightB 
+
\leftB
\begin{array}{rrr}
0 & 5 & 0 \\
4 & -4 & 14 \\
1 & 2 & 6
\end{array}
\rightB
\end{equation*}
dapat dilakukan.

Definisi formalnya adalah sebagai berikut.

\begin{definition}{Penjumlahan Matriks}\label{def:additionofmatrices}
Misalkan $A=\leftB a_{ij}\rightB $ dan $B=\leftB b_{ij}\rightB $ adalah dua
matriks $m\times n$. Maka $A+B=C$ \index{matriks!penjumlahan}, dengan $C$ sebagai matriks $m \times n$
$C=\leftB c_{ij}\rightB$ yang didefinisikan oleh
\begin{equation*}
c_{ij}=a_{ij}+b_{ij}
\end{equation*}

\end{definition}

Definisi ini menyatakan bahwa ketika menjumlahkan matriks, kita cukup menjumlahkan entri-entri matriks yang bersesuaian.
Hal ini ditunjukkan dalam contoh berikutnya. 

\begin{example}{Penjumlahan Matriks Berukuran Sama}\label{exa:samesizematrixaddition}
Jumlahkan matriks-matriks berikut jika memungkinkan.
\begin{equation*}
A = \leftB
\begin{array}{ccc}
1 & 2 & 3 \\
1 & 0 & 4
\end{array}
\rightB,
B = \leftB
\begin{array}{rrr}
5 & 2 & 3 \\
-6 & 2 & 1
\end{array}
\rightB
\end{equation*}
\end{example}

\begin{solution}
Perhatikan bahwa $A$ dan $B$ sama-sama berukuran $2 \times 3$.
Karena $A$ dan $B$ berukuran sama, penjumlahan dapat dilakukan. Dengan menggunakan Definisi \ref{def:additionofmatrices},
penjumlahan dilakukan sebagai berikut. 
\begin{equation*}
A + B = \leftB
\begin{array}{rrr}
1 & 2 & 3 \\
1 & 0 & 4
\end{array}
\rightB
+
\leftB
\begin{array}{rrr}
5 & 2 & 3 \\
-6 & 2 & 1
\end{array}
\rightB
=\allowbreak 
\leftB
\begin{array}{rrr}
1+5 & 2+2 & 3+3 \\
1+ -6 & 0+2 & 4+1
\end{array}
\rightB
=
\leftB
\begin{array}{rrr}
6 & 4 & 6 \\
-5 & 2 & 5
\end{array}
\rightB
\end{equation*}
\end{solution}

Penjumlahan matriks mengikuti sifat-sifat yang sangat mirip dengan
penjumlahan bilangan biasa. Perhatikan bahwa ketika kita menulis, misalnya, $A+B$, kita
mengasumsikan kedua matriks berukuran sama sehingga operasi tersebut
memang dapat dilakukan.

\begin{proposition}{Sifat-Sifat Penjumlahan Matriks}\label{prop:propertiesofaddition}
Misalkan $A,B$, dan $C$ adalah matriks. Maka sifat-sifat berikut \index{matriks!sifat penjumlahan} berlaku. 

\begin{itemize}
\item Hukum Komutatif Penjumlahan
\begin{equation}
A+B=B+A  \label{mat1}
\end{equation}

\item Hukum Asosiatif Penjumlahan
\begin{equation}
\left( A+B\right) +C=A+\left( B+C\right) \label{mat2}
\end{equation}

\item Keberadaan Identitas Aditif
\begin{equation}
\begin{array}{c}
\mbox{Terdapat matriks nol 0 sedemikian sehingga}\\
A+0=A  \label{mat3}
\end{array}
\end{equation}

\item Keberadaan Invers Aditif
\begin{equation}
\begin{array}{c}
\mbox{Terdapat matriks $-A$ sedemikian sehingga} \\
A+\left( -A\right) =0 \label{mat4}
\end{array}
\end{equation}
\end{itemize}
\end{proposition}

\ifdefined\showproofs
\begin{proof}
Perhatikan Hukum Komutatif Penjumlahan yang diberikan dalam \ref{mat1}. Misalkan $A,B,C,$ dan $D$ adalah matriks sedemikian sehingga $A+B=C$ dan
$B+A=D.$ Kita ingin menunjukkan bahwa $D=C$. Untuk melakukannya, kita akan menggunakan definisi penjumlahan matriks yang diberikan dalam Definisi \ref{def:additionofmatrices}.
Sekarang,
\begin{equation*}
c_{ij}=a_{ij}+b_{ij}=b_{ij}+a_{ij}=d_{ij}
\end{equation*}
Oleh karena itu, $C=D$ karena entri pada posisi $(i,j)$ sama untuk semua $i$ dan $j$. Perhatikan bahwa
kesimpulan ini mengikuti hukum komutatif penjumlahan bilangan, yang menyatakan bahwa jika $a$ dan $b$ adalah dua bilangan,
maka $a+b = b+a$.
Bukti untuk hasil-hasil lainnya serupa dan diserahkan sebagai latihan.
\end{proof}
\fi

Kita menyebut matriks nol dalam \ref{mat3} sebagai \textbf{identitas aditif}. Demikian pula, matriks $-A$
dalam \ref{mat4} disebut \textbf{invers aditif}. $-A$
didefinisikan sama dengan $\left( -1\right) A = [-a_{ij}].$ Dengan kata lain, setiap entri $A$ dikalikan dengan $-1$.
Pada bagian berikutnya, kita akan mempelajari perkalian skalar secara lebih mendalam
untuk memahami makna $\left( -1\right) A.$
